\documentclass[11pt]{article}
% Bien dich: pdflatex (hoac Overleaf). Hinh ve bang TikZ/pgfplots, khong can anh ngoai.
\usepackage[T5]{fontenc}
\usepackage[utf8]{inputenc}
\usepackage[vietnamese]{babel}
\usepackage[a4paper,margin=2.3cm]{geometry}
\usepackage{amsmath,amssymb}
\usepackage{booktabs}
\usepackage{array}
\usepackage{float}
\usepackage{caption}
\usepackage{xcolor}
\usepackage{tikz}
\usetikzlibrary{positioning,arrows.meta,calc}
\usepackage{pgfplots}
\pgfplotsset{compat=1.17}
\usepackage[hidelinks]{hyperref}

\definecolor{cOrig}{HTML}{7F7F7F}
\definecolor{cLBR}{HTML}{1F77B4}
\definecolor{cPos}{HTML}{2CA02C}
\definecolor{cNeg}{HTML}{D62728}
\definecolor{cBox}{HTML}{EAF2FB}
\definecolor{cNew}{HTML}{E8F5E9}

\newcommand{\R}{\mathbb{R}}
\newcommand{\bmu}{\boldsymbol{\mu}}
\newcommand{\bSigma}{\boldsymbol{\Sigma}}
\newcommand{\vf}{\mathbf{f}}
\newcommand{\vx}{\mathbf{x}}
\newcommand{\vz}{\mathbf{z}}
\newcommand{\vw}{\mathbf{w}}
\newcommand{\vu}{\mathbf{u}}
\newcommand{\LN}{\mathrm{LN}}
\newcommand{\muz}{\bmu^{z}}
\newcommand{\Sz}{\bSigma^{z}}

\title{Báo cáo dự án khoa học\\[4pt]
\large LBR: phát lại lớp cũ qua khối cuối của ViT\\
cho học tăng dần theo lớp không lưu ảnh với mô hình tiền huấn luyện}
\author{Thị Thanh Nga Hoàng}
\date{07/11/2026}

\begin{document}
\maketitle

\begin{abstract}
Trong học tăng dần theo lớp không lưu ảnh, các phương pháp mạnh nhất dựa trên ViT tiền huấn luyện (SLCA, SSIAT, RSIAT)
``nhớ'' mỗi lớp cũ bằng một phân phối Gauss của \emph{đặc trưng cuối} $\vf$, rồi huấn luyện lại bộ phân loại trên các
điểm giả bốc từ phân phối đó. Cách làm này ngầm giả định rằng đặc trưng thật của lớp cũ dưới mô hình hiện tại vẫn có
dạng Gauss, và rằng độ trôi của nó có thể bù bằng một phép dịch hay phép biến đổi tuyến tính trên $\vf$. Cả hai giả
định bỏ qua việc $\vf$ là đầu ra phi tuyến của khối Transformer cuối, mà adapter của khối này thay đổi ở mọi bước học.
Báo cáo trình bày \textbf{LBR} (Last-Block Replay): lưu phân phối của lớp cũ ở trạng thái trung gian $\vz$ ngay
\emph{trước} phần phi tuyến của khối cuối, bám theo độ trôi của nó, và khi căn chỉnh bộ phân loại thì cho từng điểm giả
đi qua \emph{chính} khối cuối hiện tại. LBR không lưu ảnh, không thêm tham số huấn luyện. Trên ImageNet-A (10 bước,
3 thứ tự lớp), LBR đạt $\mathbf{73.26/65.62}$ (Avg/Last) so với $72.68/65.02$ của RSIAT gốc chạy lại cùng điều kiện
($+0.58$ Avg, $+0.59$ Last; Last cao hơn ở 3/3 seed). Trên VTAB seed 1993, LBR đạt $94.97/95.10$ so với $94.72/94.69$.
Các seed còn lại của VTAB và toàn bộ CUB đang chạy.
\end{abstract}

\tableofcontents
\newpage

% =====================================================================================================================
\section{Bài toán và ký hiệu}\label{sec:setup}

\paragraph{Mô hình.} Một ảnh $\vx$ đi qua \textbf{bộ trích đặc trưng} (ViT-B/16 \cite{vit} tiền huấn luyện trên ImageNet-21K,
cộng adapter) để thành vector $\vf(\vx)\in\R^{d}$, $d=768$. \textbf{Bộ phân loại} (head) gồm một vector trọng số
$\vw_c\in\R^d$ cho mỗi lớp $c$ và dự đoán
\begin{equation}\label{eq:pred}
\hat y(\vx)=\arg\max_c\ \cos\big(\vw_c,\vf(\vx)\big),
\end{equation}
tức chọn lớp có trọng số ``cùng hướng'' nhất với đặc trưng. Độ dài của $\vf$ không ảnh hưởng tới dự đoán, chỉ hướng của
nó.

\paragraph{Học tăng dần theo lớp, không lưu ảnh.} Dữ liệu đến theo $T$ bước $t=0,1,\dots,T-1$. Bước $t$ mang tập lớp mới
$\mathcal{C}_t$ cùng ảnh huấn luyện của chúng; $K_t=|\mathcal{C}_0|+\dots+|\mathcal{C}_{t-1}|$ là số lớp đã gặp trước
bước $t$ (``lớp cũ''). Mô hình \textbf{không được giữ ảnh} của các bước trước. Sau mỗi bước, nó phải phân loại đúng mọi
lớp đã gặp mà không biết ảnh kiểm tra thuộc bước nào. Trên ImageNet-A: 200 lớp, $T=10$, mỗi bước 20 lớp, mỗi lớp
khoảng $n_c\approx30$ ảnh huấn luyện.

\paragraph{Ký hiệu khác.} $\theta_t$: tham số adapter sau bước $t$ (backbone gốc luôn đóng băng); $t_c$: bước mà lớp
$c$ được học; $\vf_{\theta}(\vx)$: đặc trưng dưới tham số $\theta$.

\paragraph{Độ đo.} Sau mỗi bước $t$, đo độ chính xác top-1 $A_t$ trên ảnh kiểm tra của mọi lớp đã gặp. $\text{Avg} =
\frac1T\sum_t A_t$, $\text{Last}=A_{T-1}$.

% =====================================================================================================================
\section{Gốc rễ vấn đề}\label{sec:problem}

\subsection{Ba nguồn gây quên}\label{sec:three}
\begin{enumerate}
\item \textbf{Trôi biểu diễn (representation drift).} Adapter được dùng chung cho mọi lớp. Khi nó học lớp mới ở bước
      $t$, đặc trưng của \emph{mọi} ảnh đều dịch chuyển, kể cả ảnh lớp cũ: $\vf_{\theta_t}(\vx)\ne\vf_{\theta_{t-1}}(\vx)$.
      Mọi thứ đã lưu về lớp cũ trong không gian đặc trưng trở nên lỗi thời.
\item \textbf{Lệch bộ phân loại.} Ở bước $t$, head chỉ thấy ảnh lớp mới. Trọng số của lớp cũ không được cập nhật theo
      đặc trưng mới, và biên giữa lớp cũ với lớp mới không được học từ dữ liệu của cả hai phía. Kết quả: ảnh lớp cũ bị
      đoán nhầm thành lớp mới.
\item \textbf{Dữ liệu thay thế không khớp đặc trưng thật.} Để sửa (2) mà không có ảnh cũ, ta phải tạo ra dữ liệu thay
      thế cho lớp cũ. Bộ phân loại học từ dữ liệu thay thế chỉ tốt bằng mức dữ liệu đó giống đặc trưng thật của ảnh lớp
      cũ dưới mô hình \emph{hiện tại}. Đây là nút thắt mà LBR nhắm vào.
\end{enumerate}

\subsection{Cách xử lý hiện nay: phát lại thống kê ở $\vf$}\label{sec:ca}
SLCA \cite{slca}, SSIAT \cite{ssiat} và RSIAT \cite{rsiat} giải quyết (2) bằng \textbf{căn chỉnh bộ phân loại}
(classifier alignment, CA). Khi học xong lớp $c$, lưu trung bình và hiệp phương sai của đặc trưng huấn luyện:
\begin{equation}\label{eq:stats_f}
\bmu^f_c=\frac1{n_c}\sum_{i:y_i=c}\vf(\vx_i),\qquad
\bSigma^f_c=\frac1{n_c-1}\sum_{i:y_i=c}\big(\vf(\vx_i)-\bmu^f_c\big)\big(\vf(\vx_i)-\bmu^f_c\big)^{\!\top}.
\end{equation}
$\bmu^f_c$ cho biết lớp $c$ nằm ở đâu trong không gian đặc trưng, $\bSigma^f_c$ cho biết nó trải rộng theo những hướng
nào. Sau mỗi bước, bốc điểm giả $\tilde\vf\sim\mathcal N(\bmu^f_c,\bSigma^f_c)$ cho \emph{mọi} lớp đã gặp và huấn luyện
lại head trên các cặp $(\tilde\vf,c)$. Vì head giờ thấy ``dữ liệu'' của cả lớp cũ lẫn lớp mới, biên giữa chúng được học
lại. Nguồn (1) được xử lý bằng cách cập nhật $\bmu^f_c$ theo độ trôi ước lượng được (mục~\ref{sec:rsiat}).

\subsection{Chỗ hở: hai giả định ẩn}\label{sec:gap}
Cách làm trên đúng nếu hai điều sau đúng:
\begin{enumerate}
\item[(G1)] Đặc trưng thật của lớp cũ \emph{dưới adapter hiện tại} có phân phối gần với một Gauss ở $\vf$.
\item[(G2)] Độ trôi của lớp cũ bù được bằng một phép dịch (hoặc biến đổi tuyến tính) trên $\vf$.
\end{enumerate}
Nhưng $\vf$ là đầu ra của khối Transformer thứ 12 (khối cuối). Viết khối này ra (adapter song song kiểu AdaptFormer
\cite{adaptformer}, như RSIAT), với $X$ là ma trận token đầu vào khối (token CLS và 196 token ảnh):
\begin{align}
X' &= X + \mathrm{MSA}\big(\LN_1(X)\big), \label{eq:attn}\\
\vf &= \LN_{\text{cuối}}\Big(X' + \mathrm{MLP}\big(\LN_2(X')\big) + s_a\,\mathrm{Adapter}_\theta(X')\Big)_{\text{CLS}}.
\label{eq:block}
\end{align}
Các đại lượng:
\begin{itemize}
\item $\mathrm{MSA}$: tự chú ý nhiều đầu (multi-head self-attention), \emph{trộn thông tin giữa các token}. Trọng số của
      nó đóng băng.
\item $\LN(\vu)=\boldsymbol\gamma\odot\frac{\vu-\bar u\mathbf 1}{\sigma(\vu)}+\boldsymbol\beta$: LayerNorm, chuẩn hoá
      từng token về trung bình 0, độ lệch chuẩn 1 rồi co giãn theo $\boldsymbol\gamma,\boldsymbol\beta$ ($\bar u$,
      $\sigma(\vu)$: trung bình và độ lệch chuẩn của các thành phần của $\vu$).
\item $\mathrm{MLP}(\vu)=W_2\,\mathrm{GELU}(W_1\vu+b_1)+b_2$: mạng truyền thẳng 768--3072--768 của ViT, đóng băng.
\item $\mathrm{Adapter}_\theta(\vu)=W_{\uparrow}\,\mathrm{ReLU}(W_{\downarrow}\vu)$: nhánh song song thắt cổ chai
      768--64--768, \emph{là phần duy nhất được học} trong khối; $s_a=0.1$ là hệ số tỉ lệ cố định.
\item Chỉ số CLS: lấy vector của token phân loại.
\end{itemize}
Ba phép trong \eqref{eq:block} (LayerNorm, GELU, ReLU) đều phi tuyến, và $\theta$ đổi ở mọi bước. Hệ quả:
\begin{itemize}
\item \textbf{Vi phạm (G1).} LayerNorm cuối ép mọi $\vf$ về gần một ``vỏ cầu'' (chuẩn hoá độ lớn, sau đó co giãn theo
      từng chiều). Một Gauss trong $\R^{768}$ không có cấu trúc này: nó đặt xác suất cả vào những vùng mà mạng không bao
      giờ sinh ra đặc trưng.
\item \textbf{Vi phạm (G2).} Khi $\theta_{t-1}\to\theta_t$, adapter khối cuối đổi, nên ánh xạ $X'\mapsto\vf$ đổi một
      cách phi tuyến. Phần trôi do khối cuối gây ra khác nhau ở từng vùng của không gian, không bắt được bằng một phép
      dịch chung cho cả lớp.
\end{itemize}

% =====================================================================================================================
\section{Phương pháp nền: RSIAT gốc}\label{sec:rsiat}
RSIAT \cite{rsiat} là điểm xuất phát và là đối chứng chính. Nó dùng một adapter cho mọi bước (số tham số không tăng).
Mỗi bước gồm bốn phần.

\paragraph{(1) Hàm mất phân loại CosFace \cite{cosface}.} Với ảnh $\vx$ nhãn $y$ và $\cos\phi_j=\cos(\vw_j,\vf(\vx))$:
\begin{equation}\label{eq:cosface}
\mathcal L_{\text{cos}}=-\log\frac{e^{\sigma(\cos\phi_y-m)}}{e^{\sigma(\cos\phi_y-m)}+\sum_{j\ne y}e^{\sigma\cos\phi_j}}.
\end{equation}
$\sigma$ (thang, 30 trên ImageNet-A) khuếch đại cosine thành logit; $m$ (biên, 0.3) bắt lớp đúng phải hơn các lớp khác
một khoảng $m$ về cosine, làm các lớp tách nhau rõ hơn so với softmax thường. Ở bước $t\ge1$ tổng chỉ lấy trên lớp mới.

\paragraph{(2) Ràng buộc biểu diễn.}
\begin{itemize}
\item Bước 0: $\mathcal L_{RS}=\frac1{|P|}\sum_{(i,j)\in P}\big[1-\cos(\vf_i,\vf_j)\big]_+ +
      \alpha\frac1{|N|}\sum_{(i,j)\in N}\big[\cos(\vf_i,\vf_j)-m_{rs}\big]_+$, với $P$ là các cặp ảnh cùng lớp, $N$ các
      cặp khác lớp trong batch, $[\cdot]_+=\max(\cdot,0)$. Số hạng đầu kéo ảnh cùng lớp về cùng hướng; số hạng sau phạt
      cặp khác lớp có cosine vượt $m_{rs}$. (Trên ImageNet-A $m_{rs}=1.5>1$, nên số hạng sau luôn bằng 0.)
\item Bước $t\ge1$: bộ tự mã hoá dư $P(\vu)=\vu+\mathrm{Dec}(\mathrm{Enc}(\vu))$ (RAE, mã 128 chiều, đầu ra sigmoid)
      được học cùng adapter, và
      \begin{equation}\label{eq:rae}
      \mathcal L_{RAE}=\beta\cdot\frac1d\big\|\vf_{\theta_t}(\vx)-P\big(\vf_{\theta_{t-1}}(\vx)\big)\big\|^2 .
      \end{equation}
      $\vf_{\theta_{t-1}}$ là đặc trưng qua bản sao đóng băng của mô hình cuối bước trước. Hàm mất buộc đặc trưng mới
      chỉ được khác đặc trưng cũ một phần dư $P(\cdot)-\mathrm{id}$ bị chặn (đầu ra sigmoid nằm trong $(0,1)$), nhờ đó
      hạn chế độ trôi mà không đóng băng hẳn. $\beta=3$ trên ImageNet-A.
\item Thêm $\gamma\,\mathcal L_{orth}$: trung bình cosine giữa $P(\bmu^f_c)$ của các lớp cũ và $P(\vf_{\theta_{t-1}}(\vx))$
      của ảnh mới, đẩy ảnh mới ra xa hướng của lớp cũ ($\gamma=6$).
\end{itemize}

\paragraph{(3) Bù trôi bằng SDC \cite{sdc}.} Với $\Delta_i=\vf_{\theta_t}(\vx_i)-\vf_{\theta_{t-1}}(\vx_i)$ trên ảnh lớp
mới:
\begin{equation}\label{eq:sdc}
\bmu^f_c\leftarrow\bmu^f_c+\frac{\sum_i \omega_{ic}\Delta_i}{\sum_i\omega_{ic}},\qquad
\omega_{ic}=\exp\!\Big(-\frac{\|\vf_{\theta_{t-1}}(\vx_i)-\bmu^f_c\|^2}{2\sigma_{sdc}^2}\Big),\ \sigma_{sdc}=4.
\end{equation}
Lớp cũ được dịch theo trung bình độ dịch chuyển của các ảnh mới nằm gần nó. Hiệp phương sai $\bSigma^f_c$ \emph{không}
được cập nhật.

\paragraph{(4) CA ở $\vf$.} Mỗi epoch, bốc 256 điểm mỗi lớp từ $\mathcal N(\kappa_c\bmu^f_c,\bSigma^f_c)$, huấn luyện
lại head cosine 10 epoch. Hệ số
\begin{equation}\label{eq:kappa}
\kappa_c=0.9+0.1\,\frac{t_c+1}{t+1}
\end{equation}
bằng 1 với lớp của bước hiện tại và nhỏ dần cho lớp học càng sớm (lớp của bước 0 ở bước 9: $\kappa=0.91$). Nhân trung
bình với $\kappa_c<1$ kéo đám mây về phía gốc toạ độ trong khi giữ nguyên độ trải; sau khi chuẩn hoá (head dùng cosine),
đám mây trở nên \emph{phân tán hơn về hướng}. Ý nghĩa: lớp càng cũ thì thống kê của nó đã qua càng nhiều lần bù trôi,
càng kém chắc chắn, nên được biểu diễn rộng hơn.

\medskip\noindent Mọi bộ nhớ và cả bù trôi lẫn CA đều ở $\vf$, nên chịu đúng hai giả định (G1), (G2).

% =====================================================================================================================
\section{Phương pháp LBR}\label{sec:lbr}

\subsection{Quan sát làm nên ý tưởng}\label{sec:obs}
Trong \eqref{eq:block}, mọi phép sau MSA (cộng dư, $\LN_2$, MLP, adapter, $\LN_{\text{cuối}}$) đều tác động \emph{riêng
lên từng token}: hàng CLS của đầu ra chỉ phụ thuộc vào hàng CLS của $X'$. Đặt
\begin{equation}\label{eq:z}
\vz:=X'_{\text{CLS}}\in\R^{768}
\end{equation}
(trạng thái của token CLS ngay sau nhánh chú ý của khối cuối). Khi đó $\vf$ là một hàm \emph{tường minh} của $\vz$:
\begin{equation}\label{eq:h}
\vf=h_\theta(\vz):=\LN_{\text{cuối}}\Big(\vz+\mathrm{MLP}\big(\LN_2(\vz)\big)+s_a\,\mathrm{Adapter}_\theta(\vz)\Big).
\end{equation}
Tính $h_\theta$ không cần ảnh, không cần 196 token ảnh, chỉ cần một vector 768 chiều. Kiểm tra số học: $\vf$ tính lại từ
$\vz$ khớp $\vf$ của mạng với sai số cỡ $6\cdot10^{-6}$.

$\vz$ là \textbf{tầng sớm nhất} có tính chất này: lùi thêm một bước (vào trong MSA) thì token CLS đã phụ thuộc vào mọi
token ảnh, không còn tái tạo được từ một vector. Đồng thời $\vz$ nằm \textbf{trước} toàn bộ phần phi tuyến có tham số
học được của khối cuối. Vì vậy:
\begin{itemize}
\item Nếu lưu phân phối của lớp cũ ở $\vz$ và cho điểm giả đi qua $h_{\theta_t}$ hiện tại, thì mọi thay đổi của adapter
      khối cuối được phản ánh \emph{chính xác}, không cần ước lượng (giảm phụ thuộc vào G2).
\item Đặc trưng giả $h_{\theta_t}(\vz)$ là đầu ra thật của mạng, nên tự động mang cấu trúc của LayerNorm và các phép
      phi tuyến (không cần G1 ở $\vf$).
\end{itemize}
Giả định Gauss vẫn còn, nhưng được đặt ở $\vz$, một tầng mà mạng chưa ``bẻ cong'' bằng phần phi tuyến cuối.

\begin{figure}[H]
\centering
\begin{tikzpicture}[node distance=6mm and 5mm,
  blk/.style={draw,rounded corners,minimum height=8mm,align=center,font=\small,fill=cBox},
  mem/.style={draw,rounded corners,minimum height=8mm,align=center,font=\small,fill=cNew},
  ar/.style={-{Latex[length=2mm]},thick}]
\node[blk] (img) {Ảnh $\vx$};
\node[blk,right=of img] (b1) {Khối 1..11\\(có adapter)};
\node[blk,right=of b1] (att) {$\LN_1$ + MSA\\khối 12};
\node[mem,right=of att,minimum width=12mm] (xp) {$\vz$};
\node[blk,right=of xp] (mlp) {$h_\theta$: MLP, Adapter,\\LN (từng token)};
\node[blk,right=of mlp] (f) {$\vf$};
\draw[ar] (img)--(b1); \draw[ar] (b1)--(att); \draw[ar] (att)--(xp); \draw[ar] (xp)--(mlp); \draw[ar] (mlp)--(f);
\node[font=\footnotesize,align=center,below=9mm of xp,text=cLBR] (lm) {LBR lưu Gauss ở đây\\$\mathcal{N}(\muz_c,\Sz_c)$};
\node[font=\footnotesize,align=center,below=9mm of f,text=cNeg] (om) {SLCA / SSIAT / RSIAT\\lưu Gauss ở đây};
\draw[cLBR,dashed,->] (lm)--(xp); \draw[cNeg,dashed,->] (om)--(f);
\end{tikzpicture}
\caption{Vị trí bộ nhớ. Đoạn từ $\vz$ tới $\vf$ chỉ gồm phép tính trên một token, nên chạy lại được trên điểm giả.}
\label{fig:where}
\end{figure}

\subsection{Tổng quan một bước $t\ge1$}
Bộ nhớ của LBR sau bước $t-1$: với mỗi lớp đã gặp, một cặp $(\muz_c,\Sz_c)$ ở tầng $\vz$; cộng head
$\{\vw_c\}$ và RAE. Một bước gồm bốn phần (Hình~\ref{fig:flow}):

\begin{center}
\fbox{\begin{minipage}{0.94\linewidth}
\textbf{Bước $t$ của LBR} (lớp mới $\mathcal C_t$, ảnh huấn luyện của lớp mới; bản sao đóng băng $\theta_{t-1}$)
\begin{enumerate}
\item \textbf{Huấn luyện adapter} trên ảnh lớp mới: $\mathcal L=\mathcal L_{\text{cos}}+\mathcal L_{RAE}$
      (mục~\ref{sec:step1}).
\item \textbf{Bám trôi ở $\vz$}: cho ảnh lớp mới qua $\theta_{t-1}$ và $\theta_t$, ước lượng phép vận chuyển, dời
      $(\muz_c,\Sz_c)$ của mọi lớp cũ (mục~\ref{sec:step2}).
\item \textbf{Ghi nhớ lớp mới} ở $\vz$: trung bình và hiệp phương sai vay hình dạng từ lớp gần nhất
      (mục~\ref{sec:step3}).
\item \textbf{CA bằng phát lại qua khối cuối}: mỗi epoch, bốc $B$ điểm $\vz$ mỗi lớp từ $\mathcal N(\muz_c,s\Sz_c)$,
      cho qua $h_{\theta_t}$, dời theo $\kappa_c$, huấn luyện lại head (mục~\ref{sec:step4}).
\end{enumerate}
Bước 0 chỉ gồm 1 (không có RAE) và 3.
\end{minipage}}
\end{center}

\begin{figure}[H]
\centering
\begin{tikzpicture}[node distance=5mm,
  st/.style={draw,rounded corners,text width=3.1cm,minimum height=11mm,align=center,font=\small,fill=cBox},
  nw/.style={st,fill=cNew},
  ar/.style={-{Latex[length=2mm]},thick}]
\node[st] (s1) {1. Huấn luyện adapter\\CosFace + RAE};
\node[nw,right=of s1] (s2) {2. Bám trôi: dời\\$\muz_c,\Sz_c$ lớp cũ};
\node[nw,right=of s2] (s3) {3. Gauss $\vz$ cho lớp\\mới (vay hình dạng)};
\node[nw,below=8mm of s3] (s4) {4a. Bốc $\vz$ giả\\mọi lớp ($s$, $B$)};
\node[nw,left=of s4] (s5) {4b. Qua khối cuối\\hiện tại: $\tilde\vf=h_{\theta_t}(\vz)$};
\node[st,left=of s5] (s6) {4c--d. Dời theo $\kappa_c$,\\huấn luyện lại head};
\draw[ar] (s1)--(s2); \draw[ar] (s2)--(s3); \draw[ar] (s3)--(s4); \draw[ar] (s4)--(s5); \draw[ar] (s5)--(s6);
\end{tikzpicture}
\caption{Luồng của LBR trong một bước $t\ge1$. Ô xanh lá là phần khác với RSIAT gốc.}
\label{fig:flow}
\end{figure}

\subsection{Bước 1: huấn luyện adapter}\label{sec:step1}
Giống RSIAT: CosFace \eqref{eq:cosface} trên lớp mới và RAE \eqref{eq:rae} ở bước $t\ge1$, cùng siêu tham số
($\sigma=30$, $m=0.3$, $\beta=3$ trên ImageNet-A; 20 epoch ở bước 0, 15 epoch ở các bước sau, SGD, batch 48). LBR đặt
$\lambda_{RS}=0$ và $\gamma=0$, tức bỏ $\mathcal L_{RS}$ và $\mathcal L_{orth}$. $\mathcal L_{orth}$ cần prototype
của lớp cũ ở $\vf$, trong khi LBR không cần lưu thống kê nào ở $\vf$.

Điểm giả của lớp cũ \emph{không} tham gia huấn luyện adapter. LBR chỉ thay đổi những gì xảy ra \emph{sau} khi adapter
học xong: cách theo dõi lớp cũ (bước 2--3) và dữ liệu cho CA (bước 4).

\subsection{Bước 2: bám trôi ở tầng $\vz$}\label{sec:step2}
\paragraph{Vì sao cần.} Phát lại qua khối cuối lo được phần trôi của khối 12. Nhưng adapter ở khối 1..11 cũng đổi (dù
có RAE neo), nên chính $\vz$ của ảnh lớp cũ cũng dịch chuyển. Bước này dời $(\muz_c,\Sz_c)$ theo phần trôi đó.

\paragraph{Dữ liệu.} Chỉ có ảnh lớp mới (không có ảnh cũ). Với mỗi ảnh huấn luyện lớp mới $\vx_i$, $i=1..N$ (biến đổi
kiểm tra, không tăng cường), lấy cặp
\[
\vz_i=\vz_{\theta_{t-1}}(\vx_i)\ \text{(mô hình đầu bước)},\qquad
\hat\vz_i=\vz_{\theta_t}(\vx_i)\ \text{(mô hình cuối bước)},\qquad \Delta_i=\hat\vz_i-\vz_i .
\]
Ý tưởng: ảnh mới đi qua cùng hai mô hình như ảnh cũ, nên cách chúng dịch chuyển cho biết không gian bị biến dạng thế
nào. Ước lượng được thực hiện một lần, sau khi huấn luyện xong bước $t$.

\paragraph{(a) Gán mềm ảnh mới cho lớp cũ.} Ảnh mới nào ``nằm trong vùng'' của lớp cũ $c$ thì độ dịch chuyển của nó
đáng tin hơn cho lớp $c$. Dùng mô hình hỗn hợp gồm $K_t$ Gauss của lớp cũ (dùng chung một hiệp phương sai) và một
thành phần nền:
\begin{equation}\label{eq:estep}
r_{ic}=\frac{\pi_c\,\mathcal N(\vz_i;\muz_c,\bar\Sigma)}
{\sum_{j=1}^{K_t}\pi_j\,\mathcal N(\vz_i;\muz_j,\bar\Sigma)+\pi_{bg}\,\mathcal N(\vz_i;\bmu_{bg},\Sigma_{bg})},
\qquad \pi_c=\frac{1-\pi_{bg}}{K_t}.
\end{equation}
\begin{itemize}
\item $r_{ic}\in[0,1]$: xác suất ảnh mới $i$ thuộc vùng của lớp cũ $c$ (``trách nhiệm'', như bước E của EM).
\item $\bar\Sigma=\Sigma_0+\frac{\operatorname{tr}\Sigma_0}{d}I$, với $\Sigma_0=\frac1{K_t}\sum_c\Sz_c$: hiệp phương sai
      chung, cộng một lượng ridge bằng phương sai trung bình mỗi chiều để ổn định khi nghịch đảo trong 768 chiều.
\item Thành phần nền $\mathcal N(\bmu_{bg},\Sigma_{bg})$: Gauss khớp trên chính các $\vz_i$ của ảnh mới (cùng cách
      ridge), xác suất tiên nghiệm $\pi_{bg}=0.5$. Nó ``hút'' những ảnh mới không gần lớp cũ nào, để chúng không làm
      nhiễu ước lượng của lớp cũ.
\item $N_c=\sum_i r_{ic}$: lượng bằng chứng (số ảnh mới hiệu dụng) cho lớp $c$.
\end{itemize}

\paragraph{(b) Phép vận chuyển chung: quay + tịnh tiến.} Mô hình hoá độ dịch chuyển của toàn bộ ảnh mới bằng
\begin{equation}\label{eq:global}
\hat\vz\approx Q(\vz-\mathbf m)+\mathbf m+\mathbf m_\Delta,\qquad
\mathbf m=\frac1N\sum_i\vz_i,\ \ \mathbf m_\Delta=\frac1N\sum_i\Delta_i ,
\end{equation}
với $Q$ là ma trận \emph{trực giao} ($Q^\top Q=I$): một phép quay (có thể kèm phản xạ) quanh tâm $\mathbf m$, cộng một
phép tịnh tiến $\mathbf m_\Delta$. Phép quay giữ nguyên khoảng cách và hình dạng, nên không thể ``bóp méo'' đám mây của
lớp cũ, chỉ xoay nó cho khớp với hệ toạ độ mới. $Q$ là nghiệm của
\begin{equation}\label{eq:procrustes}
Q=\arg\min_{Q^\top Q=I}\ \sum_{i}\big\|(\hat\vz_i-\mathbf m-\mathbf m_\Delta)-Q(\vz_i-\mathbf m)\big\|^2
+\lambda\|Q-I\|_F^2 ,
\end{equation}
bài toán Procrustes trực giao có lời giải đóng: với $M=\sum_i(\hat\vz_i-\mathbf m-\mathbf m_\Delta)(\vz_i-\mathbf m)^\top
+\lambda I=U S V^\top$ (SVD), $Q=UV^\top$. Số hạng $\lambda\|Q-I\|_F^2$ kéo $Q$ về ma trận đơn vị (``không quay'') ở các
hướng có ít dữ liệu: với $N\approx600$ ảnh trong 768 chiều, nhiều hướng không có thông tin. $\lambda=\alpha\cdot\frac1d
\sum_i\|\vz_i-\mathbf m\|^2$, với $\alpha\in\{0.01,0.1,1,10,100\}$ chọn bằng chia đôi: khớp trên nửa số ảnh của mỗi lớp
mới, đo sai số dự đoán $\Delta_i$ trên nửa còn lại, lấy $\alpha$ có sai số nhỏ nhất.

\paragraph{(c) Dời trung bình của từng lớp cũ.} Theo phép chung \eqref{eq:global}, trung bình lớp $c$ dịch một đoạn
\begin{equation}\label{eq:shift_g}
\boldsymbol\delta^{g}_c=\mathbf m_\Delta+(Q-I)(\muz_c-\mathbf m).
\end{equation}
Nếu lớp $c$ có đủ bằng chứng riêng ($N_c\ge5$), tinh chỉnh bằng các ảnh mới nằm trong vùng của nó:
\begin{equation}\label{eq:shift_l}
\bar\vz_c=\frac{\sum_i r_{ic}\vz_i}{N_c},\quad \bar\Delta_c=\frac{\sum_i r_{ic}\Delta_i}{N_c},\quad
\boldsymbol\delta^{loc}_c=\bar\Delta_c+(Q-I)(\muz_c-\bar\vz_c),\quad
\boldsymbol\delta_c=\frac{n^{\text{eff}}_c\,\boldsymbol\delta^{loc}_c+n_0^{g}\,\boldsymbol\delta^{g}_c}{n^{\text{eff}}_c+n_0^{g}},
\end{equation}
với $n^{\text{eff}}_c=(\sum_ir_{ic})^2/\sum_ir_{ic}^2$ là số ảnh hiệu dụng (bằng số ảnh nếu mọi trọng số bằng nhau, nhỏ
hơn nếu trọng số dồn vào vài ảnh) và $n_0^{g}=50$. $\boldsymbol\delta^{loc}_c$ là độ dịch trung bình của các ảnh gần
lớp $c$, hiệu chỉnh theo khoảng cách từ tâm của chúng tới $\muz_c$. Công thức trộn là ước lượng co rút: lớp có ít bằng
chứng riêng ($n^{\text{eff}}_c\ll50$) đi theo phép chung, lớp có nhiều bằng chứng đi theo độ dịch cục bộ. Lớp có
$N_c<5$ dùng $\boldsymbol\delta_c=\boldsymbol\delta^g_c$.

\paragraph{(d) Cập nhật.}
\begin{equation}\label{eq:transport}
\muz_c\leftarrow\muz_c+\boldsymbol\delta_c,\qquad \Sz_c\leftarrow Q\,\Sz_c\,Q^\top .
\end{equation}
Hiệp phương sai được xoay cùng hệ toạ độ: nếu $\vz\sim\mathcal N(\bmu,\Sigma)$ thì $Q\vz$ có hiệp phương sai
$Q\Sigma Q^\top$. Khác với SDC \eqref{eq:sdc}, cả \emph{hướng trải} của lớp cũ được cập nhật, không chỉ vị trí.

\subsection{Bước 3: ghi nhớ lớp mới, vay hình dạng}\label{sec:step3}
Sau khi bước 2 cập nhật lớp cũ, cho ảnh huấn luyện của mỗi lớp mới $c$ qua mô hình $\theta_t$ để lấy $\vz$, tính trung
bình $\muz_c$ và hiệp phương sai mẫu $C_c$ (như \eqref{eq:stats_f} nhưng ở $\vz$).

\paragraph{Vấn đề.} $C_c$ ước lượng từ $n_c\approx30$ điểm trong 768 chiều, nên có hạng tối đa $n_c-1=29$: đám mây Gauss
$\mathcal N(\muz_c,C_c)$ là một ``tấm phẳng'' 29 chiều, phương sai bằng 0 theo 739 hướng còn lại. Ảnh kiểm tra thật của
lớp $c$ lại có thành phần ở các hướng đó.

\paragraph{Lời giải: vay hình dạng của lớp gần nhất} (theo tinh thần hiệu chỉnh phân phối bằng lớp lân cận
\cite{freelunch}):
\begin{equation}\label{eq:borrow}
P_c=\sum_{j\in\mathcal S_c}a_{cj}\,\Sigma^{\text{src}}_j,\qquad
a_{cj}=\frac{\exp\big(\tau\cos(\muz_c,\muz_j)\big)}{\sum_{k\in\mathcal S_c}\exp\big(\tau\cos(\muz_c,\muz_k)\big)},\qquad
\Sz_c=\frac{n_0P_c+(n_c-1)C_c}{n_0+n_c-1}+10^{-3}I .
\end{equation}
\begin{itemize}
\item $\mathcal S_c$: tập nguồn, gồm mọi lớp cũ (nguồn $\Sigma^{\text{src}}_j=\Sz_j$, đã bám trôi ở bước 2) và các lớp
      mới khác (nguồn $\Sigma^{\text{src}}_j=C_j$), trừ chính lớp $c$.
\item $a_{cj}$: trọng số vay, softmax theo độ giống nhau về hướng giữa các trung bình. $\tau=256$ rất lớn, nên gần như
      toàn bộ trọng số dồn vào lớp gần nhất. $P_c$ vì vậy xấp xỉ hiệp phương sai của lớp láng giềng gần nhất.
\item $n_0=20$: độ tin vào hình dạng vay, tính bằng ``số ảnh giả''. Công thức chính là ước lượng hậu nghiệm kiểu Bayes:
      $P_c$ là tiên nghiệm có giá trị như $n_0$ ảnh, $C_c$ là dữ liệu có $n_c-1$ bậc tự do. Với $n_c=30$, dữ liệu riêng
      chiếm $29/49\approx59\%$, hình dạng vay $41\%$.
\item $10^{-3}I$: một lượng nhỏ để ma trận luôn xác định dương (cần cho phân rã Cholesky ở bước 4).
\end{itemize}
Ý nghĩa: các lớp gần nhau trong không gian đặc trưng của một mạng tiền huấn luyện thường có hướng biến thiên giống nhau.
Vay hình dạng lấp các hướng mà 30 ảnh không đủ để thấy.

Ở bước 0, nguồn chỉ gồm các lớp khác của bước 0.

\subsection{Bước 4: căn chỉnh bộ phân loại bằng phát lại qua khối cuối}\label{sec:step4}
Gọi $C=K_t+|\mathcal C_t|$ là tổng số lớp, $L_c$ là phân rã Cholesky của $\Sz_c$ ($L_cL_c^\top=\Sz_c$). Head được khởi
tạo từ trọng số hiện có (lớp cũ: sau CA của bước trước; lớp mới: sau bước 1). Lặp 10 epoch:
\begin{enumerate}
\item \textbf{Bốc điểm ở $\vz$.} Với mỗi lớp $c$ (cũ lẫn mới), bốc $B$ điểm
      \begin{equation}\label{eq:sample}
      \vz_{c,b}=\muz_c+\sqrt{s}\,L_c\,\boldsymbol\epsilon_{c,b},\qquad \boldsymbol\epsilon_{c,b}\sim\mathcal N(0,I),\ b=1..B,
      \end{equation}
      tức $\vz_{c,b}\sim\mathcal N(\muz_c,s\,\Sz_c)$. $s\ge1$ nới rộng đám mây.
\item \textbf{Cho qua khối cuối hiện tại}, không tính gradient: $\tilde\vf_{c,b}=h_{\theta_t}(\vz_{c,b})$ theo
      \eqref{eq:h}. Đám mây $\{\tilde\vf_{c,b}\}_b$ \emph{không} còn là Gauss: nó có hình dạng do MLP, adapter và
      LayerNorm của mạng hiện tại tạo ra.
\item \textbf{Dời theo $\kappa_c$} \eqref{eq:kappa}: $\tilde\vf_{c,b}\leftarrow\tilde\vf_{c,b}+(\kappa_c-1)\,
      \overline{\tilde\vf}_c$, với $\overline{\tilde\vf}_c=\frac1B\sum_b\tilde\vf_{c,b}$. Phép này nhân trung bình của
      đám mây với $\kappa_c$ và giữ nguyên độ lệch của từng điểm quanh trung bình, đúng như $\mathcal
      N(\kappa_c\bmu,\Sigma)$ của RSIAT nhưng áp dụng cho đám mây không Gauss.
\item \textbf{Huấn luyện lại head.} Trộn ngẫu nhiên $C\cdot B$ điểm, chia thành $C$ lô, mỗi lô $B$ điểm, và với mỗi lô
      tối ưu
      \begin{equation}\label{eq:ca_loss}
      \mathcal L_{CA}=-\frac1B\sum_{(\tilde\vf,c)\in\text{lô}}\log\frac{\exp\big(\sigma\cos(\vw_c,\tilde\vf)\big)}
      {\sum_{j=1}^{C}\exp\big(\sigma\cos(\vw_j,\tilde\vf)\big)}
      \end{equation}
      theo $\{\vw_j\}$, bằng SGD (momentum 0.9, tốc độ học giảm theo cosine qua 10 epoch). Chỉ head được cập nhật;
      adapter giữ nguyên.
\end{enumerate}
Sau CA, dự đoán theo \eqref{eq:pred}.

\paragraph{Hai siêu tham số của CA.}
\begin{itemize}
\item $s$ (độ nới): $s=1$ dùng đúng $\Sz_c$; $s>1$ làm đám mây rộng hơn, head học biên rộng hơn quanh mỗi lớp. Vai trò
      tương tự $\kappa_c$: bù cho việc bộ nhớ chưa phản ánh đủ độ bất định của lớp.
\item $B$ (số điểm mỗi lớp mỗi epoch): gradient của $\mathcal L_{CA}$ là ước lượng Monte Carlo trên các điểm giả; $B$
      lớn hơn thì ước lượng ít nhiễu hơn, nhất là khi đám mây rộng.
\end{itemize}
Mặc định $s=1$, $B=256$ (như RSIAT). Bản chính trên ImageNet-A dùng $s=2$, $B=1024$, chọn trên seed 1993
(mục~\ref{sec:grid}).

\subsection{Chi phí}
\begin{itemize}
\item \textbf{Bộ nhớ:} mỗi lớp một vector 768 chiều và một ma trận $768\times768$ ở tầng $\vz$, cùng cỡ với thống kê mà
      RSIAT lưu ở $\vf$; pipeline không cần thống kê nào ở $\vf$. (Mã hiện tại vẫn tính thêm thống kê ở $\vf$ để chẩn
      đoán; chúng không được dùng cho CA và không ảnh hưởng kết quả.) Không lưu ảnh, không thêm tham số huấn luyện.
\item \textbf{Tính toán:} bước 2 thêm một lượt cho ảnh lớp mới qua hai mô hình và một SVD $768\times768$. Bước 4, mỗi
      epoch, đẩy $C\cdot B$ vector qua MLP và adapter của một khối (chia cụm để tiết kiệm bộ nhớ GPU), cộng $C$ phân rã
      Cholesky. Với $B=256$, CA ở bước cuối của ImageNet-A mất khoảng 7 giây; chi phí tăng tuyến tính theo $B$. Huấn
      luyện adapter vẫn chiếm phần lớn thời gian.
\end{itemize}

% =====================================================================================================================
\section{Điểm khác so với các phương pháp trước}\label{sec:diff}

\begin{table}[H]
\centering\small
\caption{So sánh theo bốn câu hỏi: lưu gì, ở đâu; bù trôi thế nào; dữ liệu cho CA tạo ra thế nào; có lưu ảnh không.}
\label{tab:diff}
\begin{tabular}{p{2.5cm}p{3.0cm}p{3.2cm}p{4.2cm}p{1.2cm}}
\toprule
Phương pháp & Bộ nhớ lớp cũ & Bù trôi & Dữ liệu cho CA & Lưu ảnh\\
\midrule
L2P / DualPrompt / CODA-P \cite{l2p,dualprompt,codaprompt} & prompt + khoá & không (backbone đóng băng) & không có CA &
không\\
ADAM / APER \cite{aper} & prototype ở $\vf$ & không cần (đóng băng sau bước đầu) & không có CA & không\\
SLCA \cite{slca} & Gauss ở $\vf$ & không & Gauss ở $\vf$ & không\\
SSIAT, RSIAT \cite{ssiat,rsiat} & Gauss ở $\vf$ & SDC: dịch trung bình ở $\vf$ & Gauss ở $\vf$ & không\\
ADC / LDC \cite{adc,ldc} & Gauss / prototype ở $\vf$ & mẫu đối kháng / bộ chiếu học được, ở $\vf$ & Gauss ở $\vf$ &
không\\
\textbf{LBR} & Gauss ở $\vz$, hình dạng vay lớp gần & quay + tịnh tiến ở $\vz$, cả trung bình \emph{và} hiệp phương
sai & \textbf{$\vz$ giả qua khối cuối hiện tại}: phi tuyến, khớp adapter hiện tại & không\\
\bottomrule
\end{tabular}
\end{table}

\begin{enumerate}
\item \textbf{Vị trí bộ nhớ.} Mọi phương pháp phát lại thống kê đã công bố đều lưu ở $\vf$. LBR lưu ở $\vz$, tầng sớm
      nhất mà từ đó vẫn tái tạo được $\vf$ chính xác chỉ từ một vector.
\item \textbf{Dữ liệu giả do mạng thật tạo ra.} $\tilde\vf=h_{\theta_t}(\vz)$ là đầu ra của chính khối cuối hiện tại,
      mang cấu trúc phi tuyến của mạng, không bị giả định là Gauss ở $\vf$.
\item \textbf{Phần trôi của khối cuối được tính chính xác.} Mọi thay đổi của adapter khối cuối phản ánh trực tiếp trong
      $\tilde\vf$. Bộ bám trôi chỉ còn phải lo phần trôi của khối 1..11 ở $\vz$; SDC, ADC, LDC phải ước lượng toàn bộ độ
      trôi ở $\vf$.
\item \textbf{Cập nhật cả hình dạng.} SDC chỉ dịch trung bình. LBR xoay cả hiệp phương sai theo hệ toạ độ mới, và vay
      hình dạng từ lớp gần nhất khi số ảnh quá ít.
\end{enumerate}
LBR là thay đổi cục bộ: không đổi kiến trúc, không thêm tham số, và có thể cắm vào mọi phương pháp adapter có bước CA dựa
trên thống kê.

% =====================================================================================================================
\section{Thực nghiệm trên ImageNet-A}\label{sec:exp}

\subsection{Thiết lập}
\begin{itemize}
\item \textbf{Dữ liệu:} ImageNet-A, 200 lớp, 10 bước $\times$ 20 lớp, khoảng 30 ảnh huấn luyện mỗi lớp.
\item \textbf{Backbone:} ViT-B/16 tiền huấn luyện ImageNet-21K, adapter song song 64 chiều ở mọi khối, như RSIAT.
\item \textbf{Seed:} 1993, 1994, 1995. Seed quyết định thứ tự lớp, nên mỗi seed là một chuỗi bài toán khác nhau.
\item \textbf{RSIAT gốc:} chạy lại bằng mã và cấu hình chính thức (\texttt{exps/adapter\_imageneta.json}), cùng máy,
      cùng seed. LBR giữ mọi siêu tham số huấn luyện của RSIAT, trừ $\lambda_{RS}=\gamma=0$.
\item \textbf{Chọn siêu tham số:} chỉ $s$ và $B$ của CA, trên seed 1993, theo Avg. Seed 1994 và 1995 không được dùng
      để chọn.
\end{itemize}

\subsection{Kết quả chính}

\begin{table}[H]
\centering\small
\caption{ImageNet-A, Avg / Last (\%). Trung bình $\pm$ độ lệch chuẩn qua 3 seed. $\Delta$: chênh lệch so với RSIAT gốc
trên cùng seed.}
\label{tab:main}
\footnotesize
\setlength{\tabcolsep}{3pt}
\begin{tabular}{lcccc}
\toprule
Phương pháp & Seed 1993 & Seed 1994 & Seed 1995 & \textbf{Trung bình}\\
\midrule
RSIAT gốc \cite{rsiat} & 74.76 / 65.77 & 72.18 / 63.99 & 71.10 / 65.31 & 72.68$\pm$1.54 / 65.02$\pm$0.75\\
LBR, mặc định ($s{=}1$, $B{=}256$) & 75.32 / 66.03 & 72.70 / 64.45 & 70.93 / 64.78 & 72.98$\pm$1.80 / 65.09$\pm$0.68\\
\textbf{LBR} ($s{=}2$, $B{=}1024$) & \textbf{75.60 / 66.43} & \textbf{73.22 / 64.98} & 70.97 / \textbf{65.44} &
\textbf{73.26$\pm$1.89 / 65.62$\pm$0.60}\\
\midrule
$\Delta$ LBR mặc định & $+0.56$ / $+0.26$ & $+0.52$ / $+0.46$ & $-0.17$ / $-0.53$ & $+0.30$ / $+0.06$\\
$\Delta$ \textbf{LBR} & $+0.84$ / $+0.66$ & $+1.03$ / $+0.99$ & $-0.13$ / $+0.13$ & $\mathbf{+0.58}$ / $\mathbf{+0.59}$\\
\bottomrule
\end{tabular}
\end{table}

\begin{itemize}
\item LBR hơn RSIAT gốc $+0.58$ Avg và $+0.59$ Last. \textbf{Last cao hơn ở cả 3/3 seed}; Avg cao hơn ở 2/3 seed, riêng
      seed 1995 thấp hơn $0.13$.
\item Trên hai seed không dùng để chọn $s,B$ (1994, 1995), Last vẫn hơn $+0.99$ và $+0.13$; Avg $+1.03$ và $-0.13$.
\item Độ lệch chuẩn giữa các seed ($\approx1.5$--$1.9$ Avg) lớn hơn mức tăng, vì mỗi thứ tự lớp có độ khó khác nhau. So
      sánh theo cặp trên cùng seed ổn định hơn: trong 27 cặp (3 seed $\times$ 9 bước có CA), LBR cao hơn ở
      \textbf{22/27}.
\item Bản mặc định chỉ hơn $+0.30$ Avg và gần bằng ở Last: nới đám mây và bốc nhiều điểm hơn là cần thiết
      (mục~\ref{sec:grid}).
\end{itemize}

\begin{figure}[H]
\centering
\begin{tikzpicture}
\begin{axis}[width=0.92\linewidth,height=6.2cm,xlabel={Bước},ylabel={Độ chính xác (\%), TB 3 seed},xtick={0,...,9},
 ymin=63,ymax=86,grid=major,grid style={black!12},legend style={at={(0.02,0.03)},anchor=south west,font=\small}]
\addplot[cOrig,thick,mark=o] coordinates {(0,83.99)(1,78.08)(2,76.71)(3,74.16)(4,72.77)(5,71.65)(6,69.87)(7,68.46)(8,66.10)(9,65.02)};
\addlegendentry{RSIAT gốc}
\addplot[cLBR,very thick,mark=triangle*] coordinates {(0,84.03)(1,78.92)(2,76.40)(3,75.74)(4,73.56)(5,72.16)(6,70.66)(7,69.02)(8,66.52)(9,65.62)};
\addlegendentry{LBR ($s{=}2$, $B{=}1024$)}
\end{axis}
\end{tikzpicture}
\caption{Độ chính xác theo bước, trung bình 3 seed. Chênh lệch LBR $-$ RSIAT gốc theo bước:
$+0.04, +0.84, -0.31, +1.58, +0.79, +0.52, +0.79, +0.56, +0.43, +0.59$. Ở bước 0, khác biệt duy nhất là LBR không dùng
$\mathcal{L}_{RS}$.}
\label{fig:curve}
\end{figure}

\begin{figure}[H]
\centering
\begin{tikzpicture}
\begin{axis}[ybar,bar width=10pt,width=0.7\linewidth,height=5cm,ylabel={LBR $-$ RSIAT gốc (điểm)},
 symbolic x coords={1993,1994,1995,TB},xtick=data,ymin=-0.4,ymax=1.2,grid=major,grid style={black!12},
 legend style={at={(0.98,0.97)},anchor=north east,font=\small},enlarge x limits=0.18]
\addplot[fill=cLBR!60] coordinates {(1993,0.84)(1994,1.03)(1995,-0.13)(TB,0.58)};
\addlegendentry{Avg}
\addplot[fill=cPos!60] coordinates {(1993,0.66)(1994,0.99)(1995,0.13)(TB,0.59)};
\addlegendentry{Last}
\end{axis}
\end{tikzpicture}
\caption{Chênh lệch theo từng seed. Last dương ở mọi seed.}
\label{fig:bar}
\end{figure}

\subsection{Lớp cũ và lớp mới ở bước cuối}
\begin{table}[H]
\centering\small
\caption{Bước 9, trung bình 3 seed: độ chính xác trên 180 lớp cũ và 20 lớp mới.}
\label{tab:oldnew}
\begin{tabular}{lccc}
\toprule
Phương pháp & Tổng (Last) & Lớp cũ & Lớp mới\\
\midrule
RSIAT gốc & 65.02 & 64.60 & 68.88\\
\textbf{LBR} ($s{=}2$, $B{=}1024$) & \textbf{65.62} & \textbf{65.27} & \textbf{68.90}\\
\bottomrule
\end{tabular}
\end{table}
Toàn bộ mức tăng ở bước cuối nằm ở \textbf{lớp cũ} ($+0.67$); lớp mới giữ nguyên. LBR giảm quên mà không đánh đổi khả
năng học lớp mới.

\subsection{Tham khảo với các phương pháp đã công bố}
\begin{table}[H]
\centering\small
\caption{ImageNet-A, 10 bước, ViT-B/16-IN21K, Avg / Last. Các dòng phía trên trích nguyên văn từ Bảng 1 của SSIAT
\cite{ssiat} (seed 1993, 1996, 1997). Hai dòng cuối là của chúng tôi (seed 1993, 1994, 1995). Seed khác nhau thì thứ tự
lớp khác nhau, nên bảng này chỉ để tham khảo; so sánh chặt chẽ là Bảng~\ref{tab:main}.}
\label{tab:published}
\begin{tabular}{llcc}
\toprule
Nhóm & Phương pháp & Avg & Last\\
\midrule
Prompt & L2P \cite{l2p} & 51.24$\pm$2.26 & 44.04$\pm$0.93\\
 & DualPrompt \cite{dualprompt} & 64.59$\pm$0.08 & 53.19$\pm$0.74\\
 & CODA-Prompt \cite{codaprompt} & 63.92$\pm$0.12 & 52.08$\pm$0.12\\
PET đa dụng & LAE \cite{lae} & 58.15$\pm$0.73 & 47.18$\pm$1.17\\
 & ADA \cite{ada} & 59.43$\pm$2.20 & 50.16$\pm$0.20\\
Đóng băng + prototype & ADAM-adapter \cite{aper} & 58.84$\pm$1.37 & 48.81$\pm$0.08\\
Phát lại thống kê & SLCA \cite{slca} & 68.88$\pm$2.31 & 61.05$\pm$0.63\\
 & SSIAT \cite{ssiat} & 70.83$\pm$1.63 & 62.43$\pm$1.63\\
\midrule
\multicolumn{4}{l}{\emph{Chúng tôi chạy (seed 1993, 1994, 1995)}}\\
 & RSIAT gốc \cite{rsiat} & 72.68$\pm$1.54 & 65.02$\pm$0.75\\
 & \textbf{LBR} & \textbf{73.26$\pm$1.89} & \textbf{65.62$\pm$0.60}\\
\bottomrule
\end{tabular}
\end{table}
Chúng tôi không trích được bảng của bài RSIAT (máy chủ CVF chặn truy cập tự động), nên dòng RSIAT là số chạy lại.

% =====================================================================================================================
\section{Phân tích: mức tăng đến từ đâu}\label{sec:ablation}

\subsection{Chỉ thay dữ liệu cho CA, trên cùng một backbone}\label{sec:offline}
Câu hỏi then chốt: phát lại qua khối cuối (bước 4) có tốt hơn Gauss ở $\vf$ không, khi mọi thứ khác giữ nguyên? Ở các
run thật, head sau CA của bước này được dùng tiếp ở bước sau, nên chênh lệch giữa hai run lẫn với nhiễu tích luỹ. Vì
vậy chúng tôi lưu adapter, head trước CA và mọi thống kê ở từng bước của run LBR seed 1993, rồi chạy lại \emph{chỉ bước
CA} với các loại dữ liệu khác nhau: cùng adapter, cùng khởi tạo head, 3 seed CA. Nhánh LBR với seed CA 0 tái tạo đúng
từng con số của run thật, xác nhận cách đo là chính xác.

\begin{table}[H]
\centering\small
\caption{Backbone cố định (seed 1993), trung bình 3 seed CA, $s=1$, $B=256$. Avg tính trên bước 1--9. ``Gauss ở
$\vf$'': thống kê ở $\vf$ được bám trôi bằng cùng thuật toán của bước 2, CA như RSIAT.}
\label{tab:offline}
\begin{tabular}{lccc}
\toprule
Dữ liệu cho CA & Avg (1--9) & Last & $\Delta$Avg so với Gauss ở $\vf$\\
\midrule
Gauss ở $\vf$ (cách của SLCA / SSIAT / RSIAT) & 73.72 & 66.05 & --\\
Bộ nhớ ở $\vz$, đẩy qua khối cuối rồi khớp lại một Gauss ở $\vf$ & 73.97 & 66.34 & $+0.25$\\
\textbf{Đẩy qua khối cuối (LBR)} & \textbf{74.16} & 66.16 & $\mathbf{+0.43}$ (từng seed $+0.29, +0.54, +0.46$)\\
Đẩy qua khối cuối, không $\kappa$ & 73.90 & 66.10 & $+0.18$\\
\bottomrule
\end{tabular}
\end{table}
\begin{itemize}
\item Chỉ thay dữ liệu cho CA, LBR hơn Gauss ở $\vf$ $+0.43$ Avg, dương ở cả 3 seed CA.
\item Phần lợi gồm hai nửa cỡ ngang nhau. \emph{Vị trí bộ nhớ}: $+0.25$, đo bằng cách đẩy qua khối cuối rồi ép lại thành
      Gauss ở $\vf$ (giữ trung bình, hiệp phương sai, bỏ hình dạng). \emph{Hình dạng không Gauss do mạng tạo ra}: thêm
      $+0.18$.
\item Bỏ $\kappa$ mất $0.26$ Avg: dời theo $\kappa$ vẫn cần thiết khi phát lại qua khối cuối.
\end{itemize}

\subsection{Bỏ từng phần, chạy thật trên 3 seed}
\begin{table}[H]
\centering\small
\caption{ImageNet-A, trung bình 3 seed. Dòng 2: LBR với bước 3--4 được thay bằng CA Gauss ở $\vf$ như RSIAT (thống kê ở
$\vf$ được bám trôi bằng thuật toán của bước 2, hiệp phương sai mẫu).}
\label{tab:abl3}
\begin{tabular}{lcccc}
\toprule
Cấu hình & Avg & Last & $\Delta$Avg & $\Delta$Last\\
\midrule
RSIAT gốc & 72.68 & 65.02 & -- & --\\
LBR, thay bước 3--4 bằng CA Gauss ở $\vf$ & 72.82 & 65.61 & $+0.14$ & $+0.59$\\
LBR mặc định ($s{=}1$, $B{=}256$) & 72.98 & 65.09 & $+0.30$ & $+0.06$\\
\textbf{LBR} ($s{=}2$, $B{=}1024$) & \textbf{73.26} & \textbf{65.62} & $\mathbf{+0.58}$ & $\mathbf{+0.59}$\\
\bottomrule
\end{tabular}
\end{table}
\begin{itemize}
\item Bỏ $\mathcal L_{RS},\mathcal L_{orth}$ và thay SDC bằng bám trôi (bước 1--2) cho gần như toàn bộ mức tăng
      \emph{Last} ($+0.59$, dương ở cả 3 seed: $+0.52, +0.59, +0.65$) nhưng chỉ $+0.14$ Avg.
\item Bộ nhớ ở $\vz$ và phát lại qua khối cuối (bước 3--4, $s{=}2$, $B{=}1024$) thêm $+0.44$ Avg, giữ nguyên Last. Phần
      lợi nằm ở bước 1--7 (lớn nhất ở bước 1, 3, 4: $+1.02$, $+1.46$, $+0.92$); ở bước 8--9 gần như bằng.
\item Hai phần bổ sung cho nhau: bám trôi giữ thống kê của lớp cũ đúng chỗ đến cuối chuỗi; phát lại qua khối cuối làm dữ
      liệu cho CA sát đặc trưng thật hơn ở từng bước.
\end{itemize}

\subsection{Vay hình dạng và $\kappa$ (chạy thật, seed 1993)}
\begin{table}[H]
\centering\small
\caption{Seed 1993, Avg / Last. Mọi dòng có cùng bước 1--2.}
\label{tab:abl1}
\begin{tabular}{lcc}
\toprule
Cấu hình & Avg & Last\\
\midrule
RSIAT gốc & 74.76 & 65.77\\
CA Gauss ở $\vf$, hiệp phương sai mẫu & 75.03 & 66.29\\
CA Gauss ở $\vf$, hiệp phương sai vay \eqref{eq:borrow} & 75.06 & 66.10\\
Đẩy qua khối cuối, hiệp phương sai mẫu (không vay) & 75.00 & 65.90\\
Đẩy qua khối cuối, vay, không $\kappa$ & 75.11 & 66.03\\
Đẩy qua khối cuối, vay, có $\kappa$ (LBR mặc định) & 75.32 & 66.03\\
Như trên, $s{=}2$, $B{=}1024$ (\textbf{LBR}) & \textbf{75.60} & \textbf{66.43}\\
\bottomrule
\end{tabular}
\end{table}
\begin{itemize}
\item \textbf{Phát lại qua khối cuối cần hình dạng vay.} Với hiệp phương sai mẫu, đám mây ở $\vz$ là tấm phẳng 29 chiều;
      đẩy nó qua khối cuối không thêm thông tin ($75.00$, không hơn CA Gauss ở $\vf$). Khi đám mây có đủ chiều nhờ vay,
      phát lại qua khối cuối hơn CA ở $\vf$ cùng hiệp phương sai vay $+0.26$ Avg ($75.32$ so với $75.06$).
\item \textbf{$\kappa$ cần thiết:} bỏ $\kappa$ mất $0.21$ Avg, cùng chiều với Bảng~\ref{tab:offline}.
\item Chênh lệch giữa các dòng của một seed cỡ $0.1$--$0.3$, xấp xỉ nhiễu của CA; vì vậy phép đo trên backbone cố định
      (Bảng~\ref{tab:offline}) là bằng chứng đáng tin hơn.
\end{itemize}

\subsection{Quét $s$ và $B$ (seed 1993)}\label{sec:grid}
\begin{table}[H]
\centering\small
\caption{LBR trên seed 1993, chạy thật, Avg / Last. RSIAT gốc: $74.76 / 65.77$.}
\label{tab:grid}
\begin{tabular}{lcc}
\toprule
& $B=256$ & $B=1024$\\
\midrule
$s=1$ & 75.32 / 66.03 & 75.40 / 66.16\\
$s=1.5$ & 75.55 / 66.29 & 75.60 / 66.36\\
$s=2$ & 75.47 / \textbf{66.56} & \textbf{75.60} / 66.43\\
\bottomrule
\end{tabular}
\end{table}
Cả 6 cấu hình đều hơn RSIAT gốc ($+0.56$ đến $+0.84$ Avg; $+0.26$ đến $+0.79$ Last), nên LBR không nhạy với lựa chọn
$s,B$. Nới đám mây ($s>1$) giúp, nhất là ở Last; đây là kết quả thực nghiệm, chúng tôi chưa có giải thích đầy đủ (phép
đo ở mục~\ref{sec:fid} cho thấy ảnh huấn luyện \emph{không} co cụm hơn ảnh kiểm tra). Tăng $B$ giúp khi đám mây rộng,
đúng như ước lượng Monte Carlo cần nhiều điểm hơn khi phương sai lớn; với $s=1$, tăng $B$ chỉ thêm $+0.08$ Avg.

\subsection{Dữ liệu giả có khớp đặc trưng thật hơn không}\label{sec:fid}
So ba đám mây giả của lớp cũ với đặc trưng thật dưới mô hình hiện tại, trên $\vf$ đã chuẩn hoá (hình học mà head cosine
dùng), bằng khoảng cách năng lượng \cite{energy}
\[
\mathcal E(P,Q)=2\,\mathbb E\|X-Y\|-\mathbb E\|X-X'\|-\mathbb E\|Y-Y'\|,\qquad X,X'\sim P,\ Y,Y'\sim Q,
\]
bằng 0 khi và chỉ khi hai phân phối trùng nhau. Trung bình qua mọi lớp cũ và bước 1--9 (seed 1993):
\begin{table}[H]
\centering\small
\caption{Khoảng cách năng lượng giữa dữ liệu giả và đặc trưng thật của lớp cũ (càng nhỏ càng khớp).}
\label{tab:fid}
\begin{tabular}{lcc}
\toprule
Dữ liệu giả & So với ảnh \emph{kiểm tra} lớp cũ & So với ảnh \emph{huấn luyện} lớp cũ\\
\midrule
Gauss ở $\vf$ & 0.2162 & \textbf{0.0320}\\
Đẩy qua khối cuối rồi khớp lại Gauss & 0.2086 & 0.0449\\
\textbf{Đẩy qua khối cuối (LBR)} & \textbf{0.2070} & 0.0431\\
\bottomrule
\end{tabular}
\end{table}
\begin{itemize}
\item \textbf{So với ảnh kiểm tra}, tức phân phối mà bộ phân loại phải đối mặt, dữ liệu giả của LBR khớp hơn Gauss ở
      $\vf$ ở \emph{cả 9/9 bước}. Mức cải thiện nhỏ (khoảng 4\%) nhưng nhất quán.
\item \textbf{So với ảnh huấn luyện}, Gauss ở $\vf$ khớp hơn. Điều này dự kiến được: Gauss ở $\vf$ ước lượng từ chính
      các ảnh đó, còn LBR vay hình dạng của lớp gần nhất nên rộng hơn.
\item Tỉ số phân tán kiểm tra/huấn luyện quanh trung bình lớp là $\approx1.01$--$1.02$ ở cả $\vz$ và $\vf$: ảnh huấn
      luyện không co cụm hơn ảnh kiểm tra ở tầng nào. Lợi ích của LBR không đến từ chỗ $\vz$ ít co cụm hơn, mà từ vị trí
      bộ nhớ và việc dùng chính khối cuối để tạo dữ liệu giả.
\end{itemize}

% =====================================================================================================================
\section{VTAB và CUB}\label{sec:more}
Quy trình giống ImageNet-A: (1) chạy LBR mặc định trên seed 1993; (2) chọn $s,B$ trên backbone cố định của run đó
(mục~\ref{sec:offline}), trung bình 3 seed CA; (3) chạy thật bản đã chọn trên seed 1993, 1994, 1995; (4) so với RSIAT
gốc cùng seed. Siêu tham số huấn luyện lấy theo cấu hình RSIAT của từng dataset.

\subsection{VTAB}
50 lớp, 5 bước $\times$ 10 lớp, mỗi bước là một miền ảnh khác. Thứ tự lớp cố định (\texttt{shuffle=false}), nên seed gần
như không đổi bài toán, chỉ đổi nhiễu huấn luyện.

\begin{table}[H]
\centering\small
\caption{VTAB, chọn $s,B$ trên backbone cố định của run LBR mặc định seed 1993, trung bình 3 seed CA, Avg (cả 5 bước) /
Last.}
\label{tab:vtab_grid}
\begin{tabular}{lcc}
\toprule
& $B=256$ & $B=1024$\\
\midrule
$s=1$ & 94.86 / 94.78 & 94.84 / 94.78\\
$s=1.5$ & 94.97 / 94.96 & 94.96 / 94.96\\
$s=2$ & \textbf{95.01 / 95.08} & 94.98 / 95.02\\
\bottomrule
\end{tabular}
\end{table}
Cùng xu hướng với ImageNet-A: $s>1$ giúp. Cấu hình được chọn: $s=2$, $B=256$.

\begin{table}[H]
\centering\small
\caption{VTAB seed 1993, độ chính xác theo bước và Avg / Last (\%).}
\label{tab:vtab}
\begin{tabular}{lccccc|c}
\toprule
Phương pháp & 0 & 1 & 2 & 3 & 4 & Avg / Last\\
\midrule
RSIAT gốc & 96.70 & 93.93 & 94.00 & 94.27 & 94.69 & 94.72 / 94.69\\
LBR mặc định ($s{=}1$, $B{=}256$) & 96.90 & 94.07 & 94.34 & 94.27 & 94.79 & 94.87 / 94.79\\
\textbf{LBR} ($s{=}2$, $B{=}256$) & 96.90 & 93.86 & 94.38 & 94.61 & 95.10 & \textbf{94.97 / 95.10}\\
\bottomrule
\end{tabular}
\end{table}
\begin{itemize}
\item LBR hơn RSIAT gốc $+0.25$ Avg và $+0.41$ Last trên seed 1993. Run thật khớp đúng dự đoán từ backbone cố định
      (seed CA 0: $93.86, 94.38, 94.61, 95.10$).
\item $s,B$ được chọn trên chính seed này, nên con số trên lạc quan. Seed 1994, 1995 đang chạy. Vì thứ tự lớp cố định,
      chúng kiểm tra độ ổn định với nhiễu huấn luyện, \emph{không} phải với thứ tự lớp mới.
\item RSIAT gốc trên seed 1994, 1995: $94.86/94.74$ và $94.75/94.80$. Tham khảo từ SSIAT \cite{ssiat} (Bảng 3, Avg / bước
      cuối): SLCA $93.82/94.33$, SSIAT $94.21/94.50$.
\end{itemize}

\subsection{CUB-200}
200 lớp, 20 bước $\times$ 10 lớp. RSIAT gốc seed 1993: $93.93/90.37$ (đường do tác giả cung cấp; 15 bước đầu khớp đúng
run của chúng tôi). Bảng 1 của SSIAT dùng 10 bước nên không so trực tiếp được.

\begin{table}[H]
\centering\footnotesize
\setlength{\tabcolsep}{3.5pt}
\caption{CUB seed 1993, LBR mặc định đang chạy: độ chính xác theo bước, 13 bước đầu (\%). Chưa đủ để kết luận.}
\label{tab:cub_partial}
\begin{tabular}{lccccccccccccc}
\toprule
Bước & 0 & 1 & 2 & 3 & 4 & 5 & 6 & 7 & 8 & 9 & 10 & 11 & 12\\
\midrule
RSIAT gốc & 100.0 & 98.79 & 98.55 & 98.02 & 97.50 & 95.24 & 93.70 & 94.23 & 93.82 & 93.27 & 93.40 & 92.58 & 91.79\\
LBR mặc định & 100.0 & 98.79 & 98.55 & 98.02 & 97.68 & 95.24 & 94.34 & 94.56 & 94.20 & 93.60 & 93.40 & 92.51 & 92.12\\
\bottomrule
\end{tabular}
\end{table}
Đến bước 12, LBR bằng hoặc cao hơn RSIAT gốc ở 12/13 bước. Kết quả đủ 20 bước, bước chọn $s,B$, và 3 seed sẽ được bổ
sung khi các run hoàn tất.

% =====================================================================================================================
\section{Giới hạn}
\begin{itemize}
\item \textbf{Mức tăng vừa phải:} $+0.58$ Avg, $+0.59$ Last trên ImageNet-A. Hướng của mức tăng nhất quán (Last 3/3
      seed, 22/27 cặp seed--bước), nhưng biên độ nhỏ hơn độ lệch giữa các thứ tự lớp.
\item \textbf{Chọn $s,B$ trên seed 1993}, vốn nằm trong trung bình 3 seed. Trên hai seed còn lại, Avg là $+1.03$ và
      $-0.13$.
\item \textbf{Nguồn của mức tăng Last} là bước 1--2 (bỏ hai hàm mất phụ, bám trôi); phát lại qua khối cuối chủ yếu tăng
      Avg.
\item \textbf{Giả định Gauss vẫn còn}, chỉ được dời từ $\vf$ về $\vz$. Phần trôi của khối 1..11 vẫn phải ước lượng
      (bước 2).
\item \textbf{Chưa đủ dataset:} VTAB mới có seed 1993, CUB đang chạy.
\end{itemize}

% =====================================================================================================================
\section{Kết luận}
Nút thắt của các phương pháp phát lại thống kê (SLCA, SSIAT, RSIAT) là bộ nhớ đặt ở đầu ra $\vf$, buộc phải giả định đặc
trưng lớp cũ là Gauss và độ trôi là tuyến tính trên $\vf$, trong khi $\vf$ là đầu ra phi tuyến của khối cuối mà adapter
thay đổi ở mọi bước. LBR dời bộ nhớ về $\vz$, tầng sớm nhất mà từ đó vẫn tái tạo được $\vf$ chính xác; bám trôi cho cả
trung bình lẫn hiệp phương sai bằng phép quay và tịnh tiến; vay hình dạng từ lớp gần nhất; và tạo dữ liệu cho CA bằng
chính khối cuối hiện tại. Không lưu ảnh và không thêm tham số, LBR đạt $73.26/65.62$ trên ImageNet-A, hơn RSIAT gốc
$+0.58$ Avg và $+0.59$ Last, Last cao hơn ở mọi seed. Trên cùng một backbone, riêng việc phát lại qua khối cuối hơn
Gauss ở $\vf$ $+0.43$ Avg ở cả 3 seed CA.

\paragraph{Mã và cấu hình.} \texttt{models/lbr.py} (\texttt{xprime}: tính $\vz$; \texttt{head\_from\_xprime}: $h_\theta$;
\texttt{train\_ca\_push}: bước 4); \texttt{models/RSIAT\_adapter.py} (\texttt{\_lbr\_class\_stats}: bước 3;
\texttt{\_ca\_xprime}: gọi bước 4); \texttt{models/sts.py} (bước 2); \texttt{models/hrc.py::fixed\_shapes} và
\texttt{models/v3stats.py} (vay hình dạng, công thức~\eqref{eq:borrow}); cấu hình trong \texttt{exps/lbr2/}; đo trên
backbone cố định: \texttt{exps/lbr2/offline.py}, \texttt{exps/lbr2/tune\_offline.py}.

\begin{thebibliography}{99}
\bibitem{vit} A. Dosovitskiy et al. An Image is Worth 16x16 Words: Transformers for Image Recognition at Scale.
  \emph{ICLR}, 2021.
\bibitem{l2p} Z. Wang et al. Learning to Prompt for Continual Learning. \emph{CVPR}, 2022.
\bibitem{dualprompt} Z. Wang et al. DualPrompt: Complementary Prompting for Rehearsal-free Continual Learning.
  \emph{ECCV}, 2022.
\bibitem{codaprompt} J. S. Smith et al. CODA-Prompt: COntinual Decomposed Attention-based Prompting for Rehearsal-Free
  Continual Learning. \emph{CVPR}, 2023.
\bibitem{adaptformer} S. Chen et al. AdaptFormer: Adapting Vision Transformers for Scalable Visual Recognition.
  \emph{NeurIPS}, 2022.
\bibitem{aper} D.-W. Zhou, Z.-W. Cai, H.-J. Ye, D.-C. Zhan, Z. Liu. Revisiting Class-Incremental Learning with
  Pre-Trained Models: Generalizability and Adaptivity are All You Need. \emph{IJCV}, 2024 (arXiv:2303.07338).
\bibitem{lae} Q. Gao et al. A Unified Continual Learning Framework with General Parameter-Efficient Tuning. \emph{ICCV},
  2023 (arXiv:2303.10070).
\bibitem{ada} B. Ermis et al. Memory Efficient Continual Learning with Transformers. \emph{NeurIPS}, 2022.
\bibitem{slca} G. Zhang, L. Wang, G. Kang, L. Chen, Y. Wei. SLCA: Slow Learner with Classifier Alignment for Continual
  Learning on a Pre-trained Model. \emph{ICCV}, 2023.
\bibitem{ssiat} Y. Tan, Q. Zhou, X. Xiang, K. Wang, Y. Wu, Y. Li. Semantically-Shifted Incremental Adapter-Tuning is A
  Continual ViTransformer. \emph{CVPR}, 2024.
\bibitem{rsiat} J. Zhao, L. Huang, X. Li, Z. An, C. Yang, Y. Wang, B. Diao, Y. Xu. Representation-Steered Incremental
  Adapter-Tuning for Class-Incremental Learning with Pre-Trained Models. \emph{CVPR}, 2026.
\bibitem{sdc} L. Yu et al. Semantic Drift Compensation for Class-Incremental Learning. \emph{CVPR}, 2020.
\bibitem{adc} D. Goswami, A. Soutif-Cormerais, Y. Liu, S. Kamath, B. Twardowski, J. van de Weijer. Resurrecting Old
  Classes with New Data for Exemplar-Free Continual Learning. \emph{CVPR}, 2024.
\bibitem{ldc} A. Gomez-Villa, D. Goswami, K. Wang, A. D. Bagdanov, B. Twardowski, J. van de Weijer. Exemplar-free
  Continual Representation Learning via Learnable Drift Compensation. \emph{ECCV}, 2024.
\bibitem{cosface} H. Wang et al. CosFace: Large Margin Cosine Loss for Deep Face Recognition. \emph{CVPR}, 2018.
\bibitem{freelunch} S. Yang, L. Liu, M. Xu. Free Lunch for Few-shot Learning: Distribution Calibration. \emph{ICLR},
  2021.
\bibitem{energy} G. J. Székely, M. L. Rizzo. Energy Statistics: A Class of Statistics Based on Distances. \emph{Journal
  of Statistical Planning and Inference}, 2013.
\end{thebibliography}
\end{document}
